\chapter{国际流动积累的学术资本对科研人员早期职业发展的影响}

第四章分析了国际流动中学术资本的积累类型及其影响机制。在此基础上，本章进一步关注这些在国际流动中形成的学术资本如何在回国后的职业发展中发挥作用。结合访谈资料，本文将职业发展具体划分为初职获得、职称晋升、项目申请和机会拓展四个方面，并据此梳理文化资本、社会资本、象征资本和心理资本在不同职业发展环节中的作用关系，形成图5-1，下文将具体阐释学术资本在职业发展的四个方面中的作用。

% TODO: 插入图片 image10.png
% 图5-1 学术资本对科研人员职业发展的作用关系图

\section{学术资本对初职获得的影响}

初职获得是指博士生完成学业后，首次进入高校或科研机构担任教职或研究职位的过程，多数受访者表示，国际学术流动对于初职获得具有积极影响。根据编码结果，留学期间积累的文化资本、社会资本、象征资本与心理资本，均在博士生获得初职的过程中发挥作用。

在文化资本方面，产出型客观化文化资本中的学术成果是影响初职获得最直接的因素。多位受访者指出，留学期间产出或积累的高质量论文直接满足了招聘机构的硬性指标。例如，A05提及，"我们这边当时一定要SSCI，当时是因为那篇质量导致能进来，入职前靠的是两篇，其中一篇是中文，一篇是英文的，就是外导合作下的20年的，是SSCI一区的，这样的情况下，我这边就能入职了"。A12也指出，留学期间在导师指导下完成的论文发表满足了所在机构的入职门槛要求。在具身化文化资本方面，留学期间形成的科研能力与国际视野也在求职中发挥了作用。A07与A17都提到留学经历通过提升自身的科研能力，增强了求职时的竞争优势。A14指出，"你的国际视野，因为我们在看待一个学科的发展，例如说我来医学信息研究所，它一定是知道当时这个研究所还没有医学信息学的硕士点跟博士点，所以对于研究所未来的这种发展，我们也是一个有国际视野的，知道国际上是怎么去办这个事情的这么一个人。"

在社会资本方面，留学期间形成的高质量学术网络被认为是促进初职获得的重要因素。A14提及，"我们能够连接到例如说清华的计算机这样的资源，以及清华计算机资源能够连接到的国际的网络，以及我自己在留学期间……对我所在的机构的connection的这部分资源……对于我个人能够找到这个教职，给了我很多很多的加分。"这些社会资本为博士生的求职过程提供了差异化的竞争优势。此外，社会资本也在一定程度上对博士生初职的选择空间发挥作用，A04提到，"为自己拓宽了求职空间，就比如有些同学他可能一起当年同时在那边做联培的，然后他毕业之后重新在那边去做个博后，就留在那边去申请教职了。"

在象征资本方面，多位受访者表示，留学经历本身作为一种身份标签，在机构的招聘门槛与评价标准中被赋予了额外的认可。A10与A12提及，初职招聘机构对于留学经历有硬性的要求。A19提到，留学身份的作用随时代变化，"那时候还是鼓励出国的，不像现在，现在好像没那么鼓励了，（那时）入职的最好要有个出国的经历。"也有一些受访者表示，留学经历不是入职的决定性因素，但在同等条件下具有优势。例如，A20指出，"大家都出去比较少的情况下你出去了，可能对你来说算在同等条件下是有一点积极作用，有一点亮点。"

在心理资本方面，留学经历对部分博士生的职业方向选择产生了关键影响。例如，A22表示，"对我来说最重要的影响就是通过在那里的学习，让我觉得我还是可以做科研的，从而走上了科研道路。如果没有那里的学习，我可能也会博士毕业，但是不一定会成为一个科学家了。"对于这部分受访者而言，留学经历通过增强科研自信心和科研热情，最终使其明确了职业方向。A14也提及，留学期间找到了自己真正热爱的研究方向，科研志趣的确立直接影响了其后续的职业选择。综上所述，留学期间积累的四类学术资本在初职获得过程中发挥着不同的作用，文化资本通过提高学术成果质量、增强科研能力与国际视野提升求职竞争力，社会资本的作用体现在形成的高质量学术网络，象征资本体现在留学经历作为初职获得的硬性要求或形成差异化优势发挥作用，心理资本影响博士生的职业方向选择，进而影响初职获得（见图5-1）。

% 图5-1 学术资本对初职获得的影响路径

\section{学术资本对职称晋升的影响}

职称晋升指科研人员从初级职称（讲师、助理研究员）向副高级（副教授、副研究员）及正高级（教授、研究员）职称的晋升过程。约半数的受访者肯定了国际学术流动对其职称晋升的积极作用，根据编码结果，文化资本与象征资本对博士生的职称晋升发挥了作用。

在文化资本方面，一些受访者指出，留学期间及之后产出的高质量论文为职称晋升提供了成果基础。例如，A14表示，"在我的这两年，因为国际流动的经验，我发表的论文会比没有出国的同志们其实要好，我认为是这样的一个关系。"也有受访者表示留学期间确定的研究方向延续到回国后的科研工作中，形成持续产出，助力职称晋升。A08提及，"我所有的这些文章都是围绕着这个课题做的，课题就是在德国选定的。"

在象征资本方面，一些受访者表示，留学经历本身作为政策性加分项或硬性条件，在职称评审中享有特殊待遇。A01指出，"我们江西省它会有一些政策，对于你有这种国外交流经历，相对来说会有不同的加分……它有这样的规定，就是说你可以提交相应的材料，然后相应材料每一条它都会有不同的打分，可能这一条会给你有一定的加分。"也有受访者提到，留学经历可以替代原本必须的基层工作经历，从而加速晋升进程。A03表示，"访问学者这一年的经历是必要的，因为我们有两种选择，要么我们去下面一些分院去驻点半年，可能就是经常去到生产基地，比如说一些果园里面去实际的指导，但是这一块因为我不擅长，我主要做基础科研了，所以我就选择了CSC资助，只要是你申请了CSC资助的访问学者出去交流一年也可以用来聘职称的，它叫做一个工作经历，基层工作经历的认定，可以用这个替换。"A09的经历则展示了留学经历作为硬性要求的情况："为我15年去走那个青年教授的时候，当时他们是要求一定要有海外经历的，所以这个是有用的。"总体而言，文化资本通过产出高质量论文及研究方向延续促进职称晋升，象征资本对职称晋升的积极作用体现在政策性加分、基层工作经历替代与作为硬性要求（见图5-2）。

% 图5-2 学术资本对职称晋升的影响路径

\section{学术资本对项目申请的影响}

项目申请指科研人员申报国家级（如国家自然科学基金、国家社会科学基金）、省部级及其他科研项目的过程。约半数的受访者肯定了国际学术流动对其项目申请的积极作用，根据编码结果，国际学术流动主要通过文化资本、社会资本与象征资本对博士生的项目申请发挥作用。

在文化资本方面，具身化文化资本与客观化文化资本分别从研究能力与学术成果两个层面助力项目申请。从具身化文化资本来看，A12提及，留学期间拓展的研究方法认知与跨学科视野提供了项目的创意来源："其实我的自然科学基金的青年项目当时是在澳洲受到的一个启发，就是听他们的讲座交流，他们也经常会请nature、science的那些编辑来，当时有个老师就讲到的遥感，是空间信息技术的一些数据和技术的在他们环境领域的应用，然后我当时就结合了最核心的一个观点，就是这种遥感空间信息技术或者不同的多元数据……然后把它拓展到我这个领域来，我觉得作为一个核心的创意点拿到的自科青年，从项目层面更有说服力一点。"A21承担的国家项目以法语国家和地区为研究对象，留学期间积累的语言能力、文化理解与学术联系，成为她获得项目的关键优势。从客观化文化资本来看，留学期间产出的学术成果作为项目申报的基础，使博士生获得了竞争力。A03提到，"把我在留学中学到的这些技能，应用到我的工作中，然后就发了一些这个方向的文章，发的文章对我后续拿一些省基金、国家基金都是有帮助的。"

在社会资本方面，留学期间建立的国际学术合作网络，提升了项目申报书在团队合作方面的竞争力。A05指出，"那肯定有用……自科本子不限页数的情况下，你有什么都得往里面写，其中有一条主要是指团队跟合作，那这一块你就可以写了，我一般会写两个团队，一个团队是我国内的团队，包括我原来的博导跟我现在同事，一个是我国外的团队，那我跟我外导、我师兄师姐、我师姐的外导都有合作。"

在象征资本方面，留学身份在部分项目申报中作为硬性条件或者科研实力证明。A07指出，"有助于大家去评很多相当于人才项目，要求有海外经历"。A22提及，"（留学经历）变成了我现在的track record，变成了我科研实力的某种证明。这种情况下，它对我在国内找工作也好，对我申请项目也好，帮助都还是非常大的。"总的来说，文化资本通过学术成果的产出和研究能力的提升，社会资本通过建立国际学术合作网络，象征资本作为硬性要求或科研实力证明，对项目申请起到促进作用（见图5-3）。

% 图5-3 学术资本对项目申请的影响路径

\section{学术资本对机会拓展的影响}

机会拓展指博士生通过留学经历获得的额外职业发展机会，主要体现在访学经历上，留学期间积累的社会资本在这一过程中发挥了直接的推动作用。

在强关系网络方面，一些受访者表示，留学期间与导师建立的深厚师生关系对获得访学机会有直接作用，例如，A18提及，"中间受邀过去做访问学者的时候，又有相关的合作。我是12年去了一次，然后14年又去了一次……法国导师那边直接邀请的，他们资助的。"在弱关系网络方面，留学期间通过学术社群建立的弱联系也发挥了作用。A21提到，"比如说后面我邀请了两次法国的老师过来做讲座，那个教授不是我的导师，他是我群里一个认识的博士的导师……我就去看他的履历，他的研究方向，看完之后我觉得是有合作关系的，我后面就邀请了他。包括我今年申请的一个去法国的访学也是去那个教授的实验室。"这体现了弱关系的桥接功能，它能够帮助博士生获取异质性信息，并且通过个人的主动性，弱关系能够向强关系转化，从而拓展职业机会。社会资本对于机会拓展的影响路径如图5-4所示。

% 图5-4 学术资本对机会拓展的影响路径

\section{本章小结}

本章旨在回答第三个研究问题，即国际流动中积累的学术资本如何作用于科研人员的职业发展。围绕这一问题，本章从初职获得、职称晋升、项目申请和机会拓展四个方面展开分析。研究发现，学术资本对职业发展的影响具有明显的类型差异。在初职获得方面，文化资本、社会资本、象征资本与心理资本均发挥了作用。客观化文化资本中的论文成果是直接的竞争指标，社会资本通过高质量学术网络提升竞争力，象征资本以留学身份在招聘中发挥门槛或加分作用，心理资本帮助部分博士生确立了学术职业方向。在职称晋升方面，文化资本与象征资本发挥了主要作用，文化资本通过论文产出与研究方向延续提供成果支撑，象征资本通过政策加分、替代基层工作经历或作为硬性条件发挥制度性作用。在项目申请方面，文化资本提供了项目创意与前期成果基础，社会资本丰富了申报书的国际团队资源，象征资本作为资格条件或科研实力证明发挥作用。在机会拓展方面，社会资本发挥了主要作用，强关系网络直接推动访学机会的获得，弱关系网络通过桥接带来新的合作机会，通过向强关系转化，拓展职业发展空间。第六章将在此基础上对全文研究结论、研究启示及不足与展望进行总结。
